\documentclass[10pt,a4paper]{amsart}
\usepackage[margin=19mm]{geometry}
\usepackage{amsmath,amssymb,mathtools}
\usepackage[scr=rsfs,cal=euler]{mathalfa}
\usepackage{xfrac}
\usepackage[unicode,hidelinks,bookmarks=false]{hyperref}
\newcommand{\ie}{i.e.\ }
\newcommand{\cf}{cf.\ }
\newcommand{\resp}{respectively\ }
\newcommand{\Subsection}{\subsection}
\providecommand{\nobreakdash}{\nobreak\hbox{-}}
\setlength{\emergencystretch}{2em}
\multlinegap0pt
\usepackage{amssymb}
\usepackage{bm}
\usepackage{mathtools}
\usepackage{tcolorbox}
\usepackage{float}
\usepackage[normalem]{ulem}
\usepackage{dsfont}
\newtheorem{proposition}{Proposition}
\newcommand{\F}{\,\overline{\kern-0.175em F\kern-0em}}
\newcommand{\I}{{I}}
\newcommand{\M}{{M}}
\newcommand{\N}{{N}}
\title{\Large\textbf{{Joint Exclusivity}}}
\author{Nawaf Mohammed}
\date{Source: arXiv:2604.17490v2 (CC BY 4.0)}
\begin{document}
\maketitle

\noindent\emph{Source and licence.} \Large\textbf{{Joint Exclusivity}}. Version arXiv:2604.17490v2 (CC BY 4.0), distributed by arXiv under the Creative Commons Attribution 4.0 International licence, http://creativecommons.org/licenses/by/4.0/, as recorded in the arXiv licence metadata for this exact version and shown on the abstract page https://arxiv.org/abs/2604.17490v2. Full licence notice: \texttt{licenses/arxiv-2604.17490-CC-BY-4.0.txt}.\par\smallskip

\noindent\textit{Editorial scope.} Counted unit: the article's proposition prop:G-JE\_trivariate (source0.tex lines 1262--1275). The article's complete original proof of it is delivered with it (source0.tex lines 1277--1363, one \texttt{proof} environment of 87 lines). No mathematics is written, repaired or re-derived: the statement and the proof are the article's own lines, in the article's own notation, and no retained line is substituted or re-flowed.  Retained uncounted as the source-local context the proof needs: lines 584--587; lines 1082--1085; lines 1107--1112.  Nothing was omitted from any retained span, so the package carries no omission record.  No retained line contains a citation, so the wrapper carries no bibliography and no bibliography entry.  No retained line contains a figure environment, an image-inclusion command or drawing code, so no image is delivered and the package declares no asset dependency.  Editorial formatting only, in addition: an AMS \texttt{amsart} preamble that restates the article's own declarations and macros used by the retained lines, namely the environments \texttt{proposition} on separate counters, together with the standard packages those lines need.  The retained environments print in this document as Proposition 1 in order of appearance, whereas the article prints the same items with the numbering of its own full document.  The complete native source of the article as distributed in the arXiv e-print for this exact version is bundled unchanged as \texttt{sources/198744/source0.tex}.
\begin{equation}
\label{eq:JE_parameters_2}
\sum_{\I\in\M}(|\I|-1)\,p_{\I}\ge \sum_{i\in\N}\F_i(0)-1.
\end{equation}

\begin{equation}
\label{eq:G-JE_parameters_1}
\sum_{\I\in \M_i}p_{\I}\le G_i^*, \qquad i\in\N,
\end{equation}

\begin{equation}
\label{eq:G-JECondition}
\sum_{i\in\N}\F_i(0)-1
\le \min\left\{\frac{n-2}{n-1}\sum_{i\in\N}G_i^*,\;
\sum_{i\in\N}G_i^*-\max_{i\in\N}G_i^*\right\}.
\end{equation}

\begin{proposition}
\label{prop:G-JE_trivariate}
Let $n=3$ and suppose that condition~\eqref{eq:G-JECondition} holds. For each
$\I\in\M$ with complementary index $j\in\N\setminus\I$, define
\begin{equation}
\label{eq:G-JE_trivariate_P}
p_{\I}(\lambda)
=\lambda\,\min_{i\in\I}G_i^*
+(1-\lambda)\,\max\left\{\sum_{k\in\N}\F_k(0)-G_j^*-1,\;0\right\},
\quad\lambda\in[0,1].
\end{equation}
Then there exists $\lambda^*\in[0,1]$ such that $\{p_\I(\lambda^*)\}_{\I\in\M}$
satisfies both~\eqref{eq:G-JE_parameters_1} and~\eqref{eq:JE_parameters_2}.
\end{proposition}

\begin{proof}
Assume~\eqref{eq:G-JECondition} holds and define, for each $\I\in\M$ with
$j\in\N\setminus\I$,
\[
U_\I=\min_{i\in\I}G_i^*,\qquad
L_\I=\max\left\{\sum_{k\in\N}\F_k(0)-G_j^*-1,\;0\right\}.
\]
Here $U_\I$ is the maximum value of $p_\I$ consistent
with~\eqref{eq:G-JE_parameters_1} when all other face masses vanish, and
$L_\I$ is the lower bound on $p_\I$ derived from the combined
constraints~\eqref{eq:G-JE_parameters_1} and~\eqref{eq:JE_parameters_2}: for
$n=3$ and $\I=\{i,k\}$ with $j\in\N\setminus\I$, constraint~\eqref{eq:G-JE_parameters_1}
for index $j$ gives $p_{\{i,j\}}+p_{\{k,j\}}\le G_j^*$; substituting
into~\eqref{eq:JE_parameters_2} yields $p_\I\ge\sum_{l\in\N}\F_l(0)-G_j^*-1$,
hence $L_\I$ as stated.

The inequality $U_\I\ge L_\I$ holds by condition~\eqref{eq:G-JECondition}: for
any $j\in\N\setminus\I$,
\[
\min_{i\in\I}G_i^*+G_j^*
\ge\sum_{l\in\N}G_l^*-\max_{l\in\N}G_l^*
\ge\sum_{l\in\N}\F_l(0)-1,
\]
where the first inequality follows because $\max_l G_l^*\ge G_l^*$ for every $l$,
and the second from~\eqref{eq:G-JECondition}; rearranging gives
$U_\I\ge\sum_{l\in\N}\F_l(0)-G_j^*-1$.

Substituting $p_\I(\lambda)=\lambda U_\I+(1-\lambda)L_\I$
into~\eqref{eq:G-JE_parameters_1} and~\eqref{eq:JE_parameters_2} and solving
for $\lambda$ yields the requirement
\begin{equation}
\label{eq:G-JE_trivariate_lambda}
\ell\;\le\;\lambda\;\le\;u_i\quad\text{for all }i\in\N,
\end{equation}
where
\[
\ell=\max\left\{
\frac{\sum_{k\in\N}\F_k(0)-1-\sum_{\I\in\M}L_{\I}}
     {\sum_{\I\in\M}(U_{\I}-L_{\I})},\;0\right\},
\qquad
u_i=\frac{G_i^*-\sum_{\I\in\M_i}L_{\I}}
         {\sum_{\I\in\M_i}(U_{\I}-L_{\I})}.
\]
We show that $0\le\ell\le u_i\le1$ for every $i\in\N$.

\smallskip
\noindent\textit{Non-negativity of $u_i$.}
For fixed $i\in\N$, the numerator $G_i^*-\sum_{\I\in\M_i}L_\I$ takes one of
the following values according to how many terms in $\sum_{\I\in\M_i}L_\I$ are
strictly positive: $G_i^*$ (both zero); $G_i^*+G_j^*-(\sum_k\F_k(0)-1)$ for
$j\in\N\setminus\{i\}$ (exactly one positive); or
$\sum_k G_k^*-2(\sum_k\F_k(0)-1)$ (both positive). All three are non-negative
by condition~\eqref{eq:G-JECondition}: the first trivially; the second since
$\sum_k\F_k(0)-1\le\sum_k G_k^*-\max_k G_k^*\le G_i^*+G_j^*$ for any
$j\ne i$; the third since $\sum_k\F_k(0)-1\le\tfrac{1}{2}\sum_k G_k^*$ (first
term of~\eqref{eq:G-JECondition} with $n=3$). Hence $u_i\ge0$.

\smallskip
\noindent\textit{Upper bound $u_i\le1$.}
The bound $u_i\le1$ is equivalent to
$G_i^*\le\sum_{\I\in\M_i}U_\I=\min(G_i^*,G_j^*)+\min(G_i^*,G_k^*)$ for
$\{j,k\}=\N\setminus\{i\}$. This sum falls below $G_i^*$ only if
$G_i^*>G_j^*+G_k^*$, a strict dominance condition that can hold for at most
one index $i$ simultaneously. Hence $\min_{i\in\N}u_i\le1$.

\smallskip
\noindent\textit{Non-emptiness: $\ell\le\min_{i\in\N}u_i$.}
Since $\sum_{\I\in\M_i}(U_\I-L_\I)\le\sum_{\I\in\M}(U_\I-L_\I)$ and the
numerator of $u_i$ is non-negative, $u_i\ge(G_i^*-\sum_{\I\in\M_i}L_\I)/
\sum_{\I\in\M}(U_\I-L_\I)$.
By definition, $L_{\N\setminus\{i\}}=\max\{\sum_k\F_k(0)-G_i^*-1,0\}\ge
\sum_k\F_k(0)-G_i^*-1$, so
$\sum_k\F_k(0)-G_i^*-1-L_{\N\setminus\{i\}}\le0$. Writing
$\sum_{\I\in\M}L_\I=\sum_{\I\in\M_i}L_\I+L_{\N\setminus\{i\}}$ and adding
the non-positive term $\sum_k\F_k(0)-G_i^*-1-L_{\N\setminus\{i\}}$ to the
numerator:
\begin{align*}
u_i
&\ge\frac{G_i^*-\sum_{\I\in\M_i}L_\I
         +\sum_k\F_k(0)-G_i^*-1-L_{\N\setminus\{i\}}}
        {\sum_{\I\in\M}(U_\I-L_\I)}
=\frac{\sum_k\F_k(0)-1-\sum_{\I\in\M}L_\I}{\sum_{\I\in\M}(U_\I-L_\I)}.
\end{align*}
Since this lower bound equals the untruncated part of $\ell$, we have
$u_i\ge\ell$ for every $i\in\N$. Therefore $\lambda^*$ exists in
$[\ell,\,\min_{i\in\N}u_i]\subseteq[0,1]$.
\end{proof}

\end{document}
